% META: {"id": "1NSI-ARCHOS-EX-012", "chapitre": "1NSI-ARCHITECTURE-OS", "type_objet": "exercice", "capacites": ["1NSI-ARCHITECTURE-OS-C2"], "parcours": 2, "competences": ["programmer", "analyser", "valider"], "duree_min": 8, "mode_creation": "ex_nihilo", "sources_inspiration": [], "fichier_tex": "chapitres/1NSI-ARCHITECTURE-OS/exercices/1NSI-ARCHOS-EX-012.tex", "corrige_tex": "chapitres/1NSI-ARCHITECTURE-OS/corriges/1NSI-ARCHOS-CO-012.tex", "status": "approved", "capacites_codes": ["C2"], "methodes": ["M1"]}
% BEGIN-VERIFY
% def executer(programme, memoire):
%     registres = {}
%     for instruction in programme:
%         op = instruction[0]
%         if op == "LOAD":
%             _, reg, adresse = instruction
%             registres[reg] = memoire[adresse]
%         elif op == "ADD":
%             _, reg_dest, reg_src = instruction
%             registres[reg_dest] = registres[reg_dest] + registres[reg_src]
%         elif op == "SUB":
%             _, reg_dest, reg_src = instruction
%             registres[reg_dest] = registres[reg_dest] - registres[reg_src]
%         elif op == "STORE":
%             _, reg, adresse = instruction
%             memoire[adresse] = registres[reg]
%         elif op == "HALT":
%             break
%     return memoire, registres
% memoire = [10, 4, 0]
% programme = [
%     ("LOAD", "R0", 0),
%     ("LOAD", "R1", 1),
%     ("SUB", "R0", "R1"),
%     ("STORE", "R0", 2),
%     ("HALT",),
% ]
% memoire_finale, registres = executer(programme, memoire)
% assert memoire_finale == [10, 4, 6]
% assert registres == {"R0": 6, "R1": 4}
% END-VERIFY
\begin{exercice}{1NSI-ARCHOS-EX-002}{1}{12}
On reprend le simulateur \lstinline{executer} du cours.
\begin{enumerate}
  \item Ajouter à \lstinline{executer} le traitement de l'instruction \lstinline{("SUB", reg_dest, reg_src)}, qui effectue \lstinline{reg_dest <- reg_dest - reg_src}.
  \item Dérouler \og à la main \fg{} (sur papier) l'exécution du programme suivant, en indiquant l'état de \lstinline{registres} après chaque instruction :
\begin{python}
memoire = [10, 4, 0]
programme = [
    ("LOAD", "R0", 0),
    ("LOAD", "R1", 1),
    ("SUB", "R0", "R1"),
    ("STORE", "R0", 2),
    ("HALT",),
]
\end{python}
  \item Vérifier ta réponse en exécutant le programme avec ta fonction modifiée.
\end{enumerate}
\end{exercice}
